\section{Theoretical Background} \label{sec:theoretical-background}
Boolean satisfiability is the problem of determining whether there exists an assignment of truth values to Boolean variables that makes a propositional formula true.
A formula in conjunctive normal form (CNF) is a conjunction of clauses, where each clause is a disjunction of literals.

\noindent The movers problem is naturally modeled as a planning problem over discrete time.
Instead of searching directly over sequences of actions, we encode the existence of a valid plan of length $T$ as a SAT formula:
\[
    \Phi_T.
\]
If $\Phi_T$ is satisfiable, then there exists a valid plan of length $T$.
If it is unsatisfiable, no plan exists within that time horizon.

\noindent This approach is useful because modern SAT solvers can efficiently handle large structured Boolean formulas.
The problem is solved by increasing the time horizon until the first satisfiable formula is found:
\[
    T = 0, 1, 2, \ldots
\]
The first satisfiable $T$ corresponds to an optimal plan with respect to the number of time steps.